\emph{(i).} For $0 < \delta < 1$ the weights $k + \delta$ are positive for $k \ge 0$ and negative for $k \le -1$, where $|k + \delta| = (-k - 1) + (1 - \delta)$. So for $\operatorname{Re}s > 1$, $\sum_{k}\operatorname{sgn}(k + \delta)|k + \delta|^{-s} = \zeta(s, \delta) - \zeta(s, 1 - \delta)$ (the factor $(\kappa/r)^{-s}$ of the weights is $1$ at $s = 0$), and by (H1) (\claimref{synthesis.imported-hurwitz-zeta-values-and-dirichlet-l-functions-as-hurwitz-sums}(ii)) its value at $s = 0$ is $(\tfrac12 - \delta) - (\tfrac12 - (1 - \delta)) = 1 - 2\delta$.

\emph{(ii).} The residues $a = 2j + 1$, $0 \le j < M$, are the odd residues modulo $2M$ in $[1, 2M]$, $\delta_{j} = a/2M$, and $\chi$ vanishes on the even ones. By (H1) (\claimref{synthesis.imported-hurwitz-zeta-values-and-dirichlet-l-functions-as-hurwitz-sums}(iii)), $L(s, \chi) = (2M)^{-s}\sum_{a}\chi(a)\zeta(s, a/2M)$, so
\[
L(0, \chi) = \sum_{a}\chi(a)\bigl(\tfrac12 - \tfrac{a}{2M}\bigr) = \tfrac12\sum_{j}\chi(2j+1)(1 - 2\delta_{j}), \qquad L(-1, \chi) = -M\sum_{a}\chi(a)B_{2}\bigl(\tfrac{a}{2M}\bigr).
\]
For the trivial character the sums are $\sum_{j}(1 - 2\delta_{j}) = M - \tfrac{1}{M}\sum_{j}(2j+1) = 0$ and the stored $\sum_{j}B_{2}(\delta_{j}) = B_{2}(\tfrac12)/M$ of clause (iii) of \claimref{synthesis.light-ray-continuum-is-the-twist-integral-of-arrays-and-window-values-are-twist-averages}. The sea energy times $\beta$ of the fibre is $\tfrac{\pi}{r}B_{2}(\delta_{j})$ by \claimref{synthesis.null-ordering-of-counted-units-is-zeta-symmetric}(ii).

\emph{(iii).} The substitution $a \mapsto 2M - a$ sends $\delta_{j}$ to $1 - \delta_{j}$ and multiplies $\chi(a)$ by $\chi(-1)$. Since $1 - 2(1 - \delta) = -(1 - 2\delta)$ and $B_{2}(1 - \delta) = B_{2}(\delta)$, the first sum is multiplied by $-\chi(-1)$ and the second by $\chi(-1)$ while both are unchanged by the relabelling: the first vanishes for $\chi(-1) = 1$ and the second for $\chi(-1) = -1$.

\emph{(iv).} $M = 2$: $\delta = \tfrac14, \tfrac34$; $B_{2}(\tfrac14) = \tfrac1{16} - \tfrac14 + \tfrac16 = -\tfrac1{48} = B_{2}(\tfrac34)$; asymmetries $1 - \tfrac12 = \tfrac12$ and $1 - \tfrac32 = -\tfrac12$. The odd residues modulo $4$ are $1, 3$, and the only character not identically $1$ has $\chi_{4}(1) = 1$, $\chi_{4}(3) = -1$. With $J^{2} = -1$ (\claimref{std-algebra.complex-numbers}(ii)): $-J\,J^{a} = -J^{a+1}$ is $1, -J, -1, J$ for $a \equiv 1, 2, 3, 0$ modulo $4$, with grade-zero parts $1, 0, -1, 0$. The asymmetry sum is $\tfrac12 - (-\tfrac12) = 1 = 2L(0, \chi_{4})$; the energy sum is $-\tfrac1{48} + \tfrac1{48} = 0$, as (iii) requires of an odd character. That the halves are those of the array of Premise (Q) version 3 is clause (iii) of the cited theorem at $r = \tfrac12$.

\emph{(v).} A character $\chi_{D}$ of conductor $|D|$ read modulo $2M$ has $L(s, \chi) = L(s, \chi_{D})\prod_{p}(1 - \chi_{D}(p)p^{-s})$ over the primes $p$ dividing $2M$ and not $D$, by (H1) (\claimref{synthesis.imported-hurwitz-zeta-values-and-dirichlet-l-functions-as-hurwitz-sums}(iv)); at $s = 0$ and $s = -1$ the factors are $1 - \chi_{D}(p)$ and $1 - p\chi_{D}(p)$. With (H2) (\claimref{synthesis.imported-class-number-formula-at-zero-for-negative-discriminants}), $2L(0, \chi_{D}) = 4h(D)/w(D)$; with (H3) (\claimref{synthesis.imported-zeta-of-a-real-quadratic-field-at-minus-one-as-a-divisor-sum}) and $\zeta(-1) = -\tfrac1{12}$ (which is (H1) at $a = 1$), $L(-1, \chi_{D}) = -\tfrac15\sum_{b}\sigma((D - b^{2})/4)$. The instances: $h = 1$ and $w = 4, 2$ for $D = -4, -8$ give $1, 2$; $h(-20) = 2$ gives $4$; for $D = -3$ ($h = 1$, $w = 6$) and $D = -11$ ($h = 1$, $w = 2$) the factor $1 - \chi_{D}(2) = 2$ gives $\tfrac43$ and $4$; for $D = -7, -15, -23$, $\chi_{D}(2) = 1$ and the factor is $0$. The divisor sums for $D = 8, 12, 5, 13$ are $5, 10, 2, 10$, so $L(-1, \chi_{D}) = -1, -2, -\tfrac25, -2$; for $D = 8, 12$ no prime is added and the energy sums are $\tfrac14$ and $\tfrac26 = \tfrac13$; for $D = 5, 13$ the factor $1 - 2\chi_{D}(2) = 3$ gives $\tfrac{3 \cdot 2/5}{5} = \tfrac{6}{25}$ and $\tfrac{3 \cdot 2}{13} = \tfrac{6}{13}$.
